\subsection{A consistent local Cartan--scalar reduction and current response}
\label{sec:cartan-scalar-reduction}

Work on the flat Spin chart and local classical action of
Theorem~\ref{thm:local-sm-jet-action}. All kinetic coefficients are positive
constants; the mass, quartic and Yukawa coefficients remain declared. Use
the standard weak Cartan generator $i\sigma_3/2$ and central generator $i$.
Writing the two real Cartan connection coordinates as $a=(a_2,a_1)$ gives
\[
 q=(1/2,1/2)^T,\qquad K=\operatorname{diag}(\kappa_2,2\kappa_1),\qquad
 d=q^TK^{-1}q,\qquad w=K^{-1}q/d.
\]
Here the gauge density is $-\tfrac14 f^TKf$: the fundamental traces are
$\operatorname{Tr}((i\sigma_3/2)^2)=-1/2$ and $\operatorname{Tr}(i^2)=-1$.
The executable weak basis uses $i\sigma_3$, in which
$q=(1,1/2)^T$ and $K=\operatorname{diag}(4\kappa_2,2\kappa_1)$; $d$ agrees.

\begin{theorem}[Local classical solution embedding]
\label{thm:cartan-scalar-reduction}
Set all fermions and their independent conjugates, color connections and
off-diagonal weak connections to zero. Let $H=(\varphi,0)$ and $a=wA$.
Every smooth solution, including the temporal gauge equation, of
\begin{equation}
 \mathcal L_{\rm red}=-\frac{F_{\mu\nu}F^{\mu\nu}}{4d}
  +| (\partial_\mu+iA_\mu)\varphi|^2
  -m^2|\varphi|^2-\lambda|\varphi|^4
 \label{eq:cartan-reduced-action}
\end{equation}
lifts to a solution of the full local classical equations, with
ansatz-compatible initial and boundary data. All complex Yukawa entries
remain unrestricted.
\end{theorem}
\begin{proof}
The identities $q^Tw=1$, $Kw=q/d$ and $w^TKw=1/d$ give the density and both
Cartan equations. In the current convention
$\delta\mathcal L_{\rm matter}=j^\mu\delta A_\mu$, these equations are
$q_a[d^{-1}\partial_\nu F^{\nu\mu}+j^\mu]=0$.
If $q^Tn=0$, the decomposition $a=wA+nB$ has no scalar coupling to $B$ and
no kinetic cross term, since $w^TKn=0$; $B=0$ solves its free equation.
The Higgs and its covariant derivative lie in the upper weak line, so both
off-diagonal weak currents vanish. The gauge covariant divergence of a
Cartan curvature is Cartan. Color curvature and sources vanish. The
diagonal connection preserves the lower Higgs line, whose potential
derivative vanishes at zero. Every fermion equation contains a fermion
factor; every Yukawa contribution to a Higgs equation contains two. Thus
all discarded variations vanish without setting any Yukawa coefficient to
zero. The temporal Cartan equation is included, so Gauss's law is retained.
\end{proof}

This is a local real-connection statement. Generic kinetic coefficients
give a direction incommensurate with the full charge lattice. No compact
subgroup, global quotient, large-gauge identification, electromagnetic
vacuum, perturbative stability or quantum truncation follows. Retaining a
real unwrapped connection is essential: its exponential alone need not
determine the embedded connections in every representation.

\paragraph{A compatible spatial discretization.}
The same reduction holds for a declared semidiscrete gauge--scalar action
with covariant link differences, invariant electric metric, and the local
logarithmic plaquette energy. For the weak holonomy use its principal
matrix logarithm on the smooth identity chart. On the Cartan slice,
$\log U_f=i\sigma_3 w_{\rm w}(\mathrm dA)_f/2$ when
$|w_{\rm w}(\mathrm dA)_f|<2\pi$, where $w_{\rm w}$ is the weak component of $w$.
The hypercharge curl is real and unwrapped. Consequently the plaquette
energy restricts to the quadratic Maxwell curl with coefficient $1/d$.
The first variation of $\operatorname{Tr}(\log U)^2$ is
$2\operatorname{Tr}(\log U\,U^{-1}\delta U)$; discarded weak variations
are off-diagonal and have zero trace pairing. Scalar link currents also
have zero off-diagonal component. Electric transport along the Cartan line
is geodesic for the invariant metric. These facts give the same omitted
equations and weighted Cartan square. This statement requires the logarithm
chart throughout the history. A Wilson cosine plaquette does not supply
this quadratic reduction for arbitrary $w$.

On the prepared golden $q=5$ tensor family use dimensionless positions
$x/L$, $L=2/\sqrt{\phi+2}$. The ordered one-dimensional nodes are
$0,2\phi-3,4\phi-6,\phi-1,3\phi-4,1$. The $4^3$ interior nodes have
lumped masses $v_i$; an internal edge has conductance $c_e$ equal to its
transverse dual area divided by its gap. The fixed zero scalar boundary
contributes the usual nonnegative onsite stiffness $b_i$. Declare
$\kappa_1=\kappa_2=m^2=1$, $\lambda=1/4$, so $d=3/8$.
For each internal plaquette set its Maxwell weight $\gamma_f$ to $1/d$
times normal dual length divided by area, and set $\epsilon_e=c_e/d$.
These are action quadratures, not causal counting weights.

With $\pi_i=v_iD_0\psi_i$ and canonical real link momentum $\Pi_e$, the
Hamiltonian is
\begin{align}
 \mathcal H={}&\sum_i\frac{|\pi_i|^2}{v_i}
   +\sum_e\frac{\Pi_e^2}{2\epsilon_e}
   +\sum_{e=i\to j}c_e|e^{ia_e}\psi_j-\psi_i|^2\notag\\
 &+\sum_i[(v_i+b_i)|\psi_i|^2+v_i|\psi_i|^4/4]
   +\tfrac12\sum_f\gamma_f(\mathrm da)_f^2 .
\end{align}
The complex scalar symplectic form is
$2\operatorname{Re}(\mathrm d\overline\pi_i\wedge\mathrm d\psi_i)$.
Differentiating this same energy gives
\[
 J_{i\to j}=-2c_e\operatorname{Im}(\overline\psi_i e^{ia_e}\psi_j),\qquad
 \rho_i=2\operatorname{Im}(\overline\psi_i\pi_i),\qquad
 G_i=\rho_i+(\operatorname{div}\Pi)_i.
\]
Here divergence is outward; $\dot\rho+\operatorname{div}J=0$ and
$\dot\Pi=J-\mathrm d^T\gamma\mathrm da$ imply $\dot G=0$.

\begin{proposition}[Local split evolution preserves Gauss constraints]
\label{prop:cartan-scalar-split-gauss}
Scalar kinetic drift, electric kinetic drift and every individual edge,
onsite and plaquette potential kick preserve each $G_i$ in ideal
arithmetic. Their symmetric composition is symplectic and preserves the
initial Gauss constraints. Disjoint local subflows commute.
\end{proposition}
\begin{proof}
Scalar drift adds a real multiple of $\pi_i$ to $\psi_i$, leaving $\rho_i$
unchanged; electric drift changes only $a$. An edge kick changes the two
charges by $-hJ_e,+hJ_e$ and $\Pi_e$ by $hJ_e$. Onsite kicks add a real
multiple of $\psi_i$ to $\pi_i$. A plaquette kick adds a curl transpose to
$\Pi$, whose divergence vanishes. Each map is its Hamiltonian subflow;
disjoint subflows have disjoint dynamical variables. The potential kicks
also commute because their coordinates are fixed during the kicks.
\end{proof}

\paragraph{Finite current and feedback control.}
The executable packet uses 64 scalar sites, 144 internal links and 108
plaquettes, with two steps of declared model duration $1/200$. It compares
the real preparation $\psi_i=1/5$, $\pi_i=a_e=\Pi_e=0$ with a local
phase intervention $(3+4i)/5$ and a genuine gauge copy, which also changes
the incident unwrapped connections. The intervention leaves charge and
Gauss constraints unchanged but changes the edge current. At the selected
edge $c_e=1/4$ exactly, so $J_e=2/125$ and its first electric half-kick is
$1/25000$. The gauge copy preserves invariant readouts. Complete operation
records bind every dynamic read to its version's actual writer; decoded
checkpoints contain scalar intensity, charge, edge current and electric
feedback. Independent high-precision replay checks these binary64 records,
including a nonempty off-influence control. This numerical agreement is
separate from the exact local identities and gives no continuous-history
error enclosure. The prepared addresses, local execution law, action,
boundary and model tick remain supplied. No causal-count clock attachment,
inherited fine-grid detector error, selected physical current or full
Standard Model quantum evolution is asserted.
